% ============================================================
%  R. Nielsen, "Den propædeutiske Logik."
%  Kjøbenhavn: Forlagt af Boghandler P. G. Philipsen;
%  Trykt i BiancoLunos Bogtrykkeri, 1845.
%  TRANSCRIPTION (Danish) — from the Royal Danish Library scan.
%
%  Scan: ~/bibliotek/Nielsen, Rasmus/1845-propædeutiske-logik.pdf
%        (KB / ABBYY Recognition Server, 297 PDF pp., A4, 105 MB. A byte-for-byte
%        twin, propædeutiske_logik.pdf, sits in the same directory — same scan.)
%  TYPE: Fraktur (Danish body); Latin/French/English and all figures in antiqua
%        -> \textit{}. Letterspacing (Sperrsatz) -> \emph{}; image-verify by zoom.
%
%  OFFSET (VERIFIED, UNIFORM): PDF = printed + 9.
%        Body printed pp. 1--283 = PDF 10--292. No leaf was scanned twice and the
%        offset never changes. Verified three ways: the ABBYY layer's own header
%        numerals (224 pages at exactly +9, the 7 exceptions all single-digit
%        misreads with +9 neighbours); a text-similarity sweep over adjacent and
%        next-but-one page pairs, which found no duplicated leaf; and the image at
%        200 dpi for the endpoints and the numeral-less opening run
%        (PDF 11=2, 12=3, 13=4, 14=5, 16=7, 292=283).
%        pagemap.py is the single source of truth — do not hard-code +9.
%
%  EXTENT: body printed pp. 1--283 (PDF 10--292). Front matter: title page PDF 6,
%        Indhold PDF 8--9. There is NO half-title, NO dedication, NO motto and
%        NO Forord. Back matter: "Rettelser" errata PDF 293; PDF 294--297 blank.
%
%  STRUCTURE (from the Indhold, PDF 8--9, transcribed from the image):
%        Indledning (I. Den formelle og speculative Logik; II. Den propædeutiske
%        Logiks Formaal) — Første Deel: Læren om det subjective Begreb
%        (Cap. 1 §§1--5, Cap. 2 §§6--8, Cap. 3 §§9--11) — Anden Deel: Læren om det
%        objective Begreb (Indledning §12, Cap. 1 §§13--15, Cap. 2 §§16--18,
%        Cap. 3 §§19--21).
%        The §§ run CONTINUOUSLY 1--21 across both Parts; they do not restart.
%        The Første Deel Indledning's two sections are numbered I./II., NOT as §§;
%        the Anden Deel Indledning IS a § (§ 12). That asymmetry is in the book.
%        The Indhold carries NO page numbers at all — bare text, no leaders.
%
%  PRINTED-vs-INDHOLD VARIANTS (record both, do NOT normalise — see CLAUDE.md):
%        body heads use a COLON  ("Første Deel:"),  Indhold a PERIOD ("Første Deel.")
%        printed p.97 has LOWERCASE "anden Deel:" against p.6's "Første Deel:"
%           and against the Indhold's "Anden Deel."
%
%  ERRATA (from "Rettelser", PDF 293; apply inline, flagged with % ERRATA):
%        p. 80  l.3  fra oven : Omsætningen   -> Oversætningen
%        p. 92  l.13 fra oven : A er ei = A,  -> A er ei = --- A
%              [the errata line itself misprints "ci" for "ei"; see RESUME-NOTES]
%        p. 93  l.7  fra oven : Betingelser   -> Betingethed
%        p.112  l.4  fra neden: Daddel        -> Dadel
%        p.152  l.8  fra neden: med           -> ved
%        p.154  l.15 fra neden: Dyr=Rige      -> Dyre=Rige
%        p.187  l.9  fra neden: εἵδη          -> εἴδη
%        p.206  l.9  fra neden: en miniatur   -> en miniature
%        p.241  l.1  fra neden: θετον         -> θεῖον   [script theta normalised]
%
%  NO FOOTNOTES anywhere in this book (confirmed on the image and by a type-size
%        sweep of all 283 pages). References are inline in parentheses; the
%        "Anm." remarks do the work footnotes would do. The "3*", "12*", "18*"
%        at some page feet are SIGNATURE MARKS, not note marks.
%  NO figures, NO tables, NO displayed equations. Inline antiqua symbols only
%        (A=A, A er ei = --- A, A/B/C placeholders). Polytonic Greek occurs inline
%        (pp. 138, 147, 179, 187, 241 ...). The fount prints theta as the script
%        form and kappa sometimes as the cursive form; both are NORMALISED to
%        θ and κ per the repo's standing convention (fount variants, not distinct
%        letters — see texts/nielsen/religionsphilosophie/transcription.tex).
%        This is also a hard build requirement: U+03D1 ϑ and U+03F0 ϰ are FATAL
%        errors under textalpha/greek-fontenc in text mode. `verify.sh` lints for
%        them, because the sandbox cannot compile Greek at all (textalpha is not
%        installed there and tlmgr cannot cross-release install it).
%  RUNNING HEADS: none in the original — only a centred bold antiqua folio.
%
%  SCAN CAVEAT: printed p.96 (PDF 105) carries a previous reader's yellow
%        highlighter across the Sibbern quotation. Not a printing feature.
%
%  STATUS: IN PROGRESS. See RESUME-NOTES.md.
% ============================================================
\documentclass[12pt]{book}
\usepackage[utf8]{inputenc}
\usepackage[T1]{fontenc}
\usepackage[danish]{babel}
\usepackage[protrusion=true,expansion=true]{microtype}
\usepackage[margin=1.4in]{geometry}
\usepackage{setspace}
\usepackage{amsmath}
\usepackage{graphicx}    % for \reflectbox (the old Danish »ɔ:« id-est mark)
\DeclareUnicodeCharacter{0254}{\reflectbox{c}}  % ɔ (U+0254) = reversed c
\usepackage{libertinus}
\usepackage{libertinust1math}
\usepackage{textalpha}   % polytonic Greek typed directly
\usepackage{fancyhdr}
% plainpages=false + pdfpagelabels: the title page, the roman front matter and the
% arabic body would otherwise all produce a PDF destination called "page.1", and
% pdfTeX drops the duplicates with a warning on every build.
% The PDF metadata is set with \hypersetup AFTER \begin{document}, NOT as package
% options here. As a package option the "æ" of the title reaches hyperref while
% fontenc still has it as the T1 command \T1\ae, which hyperref cannot put in a
% PDF string: it warns "Token not allowed in a PDF string" and drops the letter.
% By \begin{document} the encodings have settled and the æ comes through intact.
\usepackage[colorlinks=true,linkcolor=black,urlcolor=black,
  unicode=true,plainpages=false,pdfpagelabels]{hyperref}
\setlength{\headheight}{15pt}
\pagestyle{fancy}
\fancyhf{}
\fancyhead[LE,RO]{\thepage}
\fancyhead[RE]{\textit{Den propædeutiske Logik}}
\fancyhead[LO]{\textit{\leftmark}}
\setlength{\parindent}{1.5em}
\setlength{\parskip}{0pt}
\onehalfspacing

% ------------------------------------------------------------
%  STRUCTURAL MACROS — use these, do not hand-roll centre blocks.
%  They exist so that every batch renders the book's four recurring
%  display forms identically, and so that check.py can count them.
% ------------------------------------------------------------

% A Part opening. In the book this is a page-top head on an ordinary text page
% (NOT a display page), introduced by a full-measure double rule.
%   #1 = the Deel line as printed, WITH its colon ("Første Deel:")
%   #2 = the Deel argument ("Læren om det subjective Begreb.")
% VERIFIED on printed p.6 at 400 dpi (batch 1): "Den propædeutiske Logiks" and
% the Deel line are bold but NOT letterspaced; only the Deel ARGUMENT is
% letterspaced. And there is NO short centred rule between the argument and the
% Capitel head that follows — the earlier note claiming one was wrong.
\newcommand{\deel}[2]{%
  \par\vspace{1.2\baselineskip}%
  \noindent\rule{\linewidth}{1.2pt}\par\vspace{2pt}%
  \noindent\rule{\linewidth}{0.4pt}\par\vspace{0.9\baselineskip}%
  \begin{center}
    {\large\bfseries Den propædeutiske Logiks\par}\vspace{0.35\baselineskip}
    {\bfseries #1\par}\vspace{0.35\baselineskip}
    {\Large\bfseries\emph{#2}\par}
  \end{center}
  \vspace{0.5\baselineskip}
  \phantomsection
  \addcontentsline{toc}{section}{\texorpdfstring{#1\quad #2}{#1 #2}}%
  \markboth{#1\ #2}{#1\ #2}%
}

% A Capitel head: bold letterspaced label, then the argument one size down.
\newcommand{\capitel}[2]{%
  \par\vspace{0.9\baselineskip}%
  \begin{center}
    {\bfseries\emph{#1}\par}\vspace{0.25\baselineskip}
    {\large #2\par}
  \end{center}
  \vspace{0.3\baselineskip}
  \phantomsection
  \addcontentsline{toc}{subsection}{\texorpdfstring{#1\quad #2}{#1 #2}}%
}

% A § head: centred "§ N." (antiqua figures), then the centred BOLD argument.
% Mid-page § heads get extra space above and below but NO rule.
%   \parag{1}{Almeenforestillingen.}
% VERIFIED on printed p.6 (batch 1): the § argument is bold but NOT letterspaced,
% unlike the Capitel label and the Deel argument. No \emph here.
\newcommand{\parag}[2]{%
  \par\vspace{0.9\baselineskip}%
  \begin{center}
    {\bfseries \S\,#1.\par}\vspace{0.2\baselineskip}
    {\bfseries #2\par}
  \end{center}
  \vspace{0.25\baselineskip}
  \phantomsection
  \addcontentsline{toc}{subsubsection}{\texorpdfstring{\S\,#1.\quad #2}{§ #1. #2}}%
}

% The Indledning's two sections, headed with a bare roman numeral. These are
% numbered I./II. and are NOT §§ -- the asymmetry is in the book (the Anden
% Deel's Indledning IS a §, § 12). Image evidence (printed pp. 1 and 4): a
% centred bold roman numeral in antiqua, then the argument in bold Fraktur a
% size up from the body, NOT letterspaced.
%   \romsec{I.}{Den formelle og speculative Logik.}
% NB the \texorpdfstring is load-bearing: a bare \quad inside \addcontentsline
% makes hyperref warn "Token not allowed in a PDF string" on every build.
\newcommand{\romsec}[2]{%
  \par\vspace{0.9\baselineskip}%
  \begin{center}
    {\large\bfseries #1\par}\vspace{0.3\baselineskip}
    {\large\bfseries #2\par}
  \end{center}
  \vspace{0.3\baselineskip}
  \phantomsection
  \addcontentsline{toc}{subsection}{\texorpdfstring{#1\quad #2}{#1 #2}}%
}

% A top-level division head that is not a Deel, a Capitel or a §. In this book
% that is "Indledning." -- which occurs TWICE: once at printed p.1 opening the
% whole work, and again at printed p.97 opening the Anden Deel.
\newcommand{\division}[1]{%
  \par\vspace{1.0\baselineskip}%
  \phantomsection
  \addcontentsline{toc}{section}{#1}%
  \markboth{#1}{#1}%
  \begin{center}{\Large\bfseries #1}\end{center}
  \begin{center}\rule{0.35\linewidth}{0.4pt}\end{center}
  \vspace{0.3\baselineskip}%
}

% An "Anm." remark. NOT a smaller type block — measured on printed p.7, the
% baseline pitch is identical to body text. The ONLY distinction is the bold
% letterspaced lead-in, after which the remark runs on as an ordinary indented
% paragraph and simply ends. Use as:  \anm{2} Text of the remark...
\newcommand{\anm}[1]{\emph{\textbf{Anm.~#1.}}}

% The Danish "that is" sign, ɔ: — type it with this, never as a plain "o:".
\newcommand{\dsi}{\textup{ɔ:}}

\begin{document}
\hypersetup{pdftitle={Den propædeutiske Logik},pdfauthor={Rasmus Nielsen}}

% \frontmatter must come BEFORE the titlepage. The titlepage environment resets
% the page counter to 1 in whatever numbering is in force; if that is arabic, the
% title page and the first page of \mainmatter both claim the PDF destination
% "page.1" and pdfTeX drops one with a warning on every build.
\frontmatter

% ============================================================
% TITLE PAGE — PDF 6
% ============================================================
\begin{titlepage}
\centering
\vspace*{3.0cm}
{\Huge\bfseries Den propædeutiske Logik\par}
\vspace{1.6cm}
{\small af\par}
\vspace{0.7cm}
% "Lie. theol." is what the page actually prints: the third glyph carries a
% clear horizontal crossbar and is identical to the 'e' of "theol.". A wrong-sort
% error for "Lic. theol." (licentiatus theologiæ). Transcribed as printed.
{\large\textit{Lie. theol.}\ {\LARGE\bfseries R. Nielsen,}\par}
\vspace{0.4cm}
{\large Professor i Philosophien.\par}
\vspace{1.8cm}
{\Large $\ast$\par}          % the fleuron: a double-taper rule with a centre lozenge
\vspace{1.8cm}
{\Large\bfseries Kjøbenhavn.\par}
\vspace{0.6cm}
% "BiancoLunos" is set solid as one word on the title page: measured at 300 dpi
% the o|L gap is 14 px against 25--27 px for the true word spaces of that line.
{\emph{Forlagt af Boghandler P. G. Philipsen.}\par}
\vspace{0.35cm}
{\small\emph{Trykt i BiancoLunos Bogtrykkeri.}\par}
\vspace{1.2cm}
{\large\textit{1845.}\par}
\end{titlepage}


% ============================================================
% INDHOLD — PDF 8--9, transcribed from the image.
% The original carries NO page numbers against any entry: bare text, no leaders,
% nothing in the right margin. Reproduced here as printed; the generated
% \tableofcontents below carries this edition's own page numbers.
% ============================================================
\chapter*{Indhold}
\markboth{Indhold}{Indhold}
\begin{center}\rule{0.35\linewidth}{0.4pt}\end{center}
{\small\setlength{\parindent}{0pt}\setlength{\parskip}{2pt}

\begin{center}\emph{Indledning.}\end{center}
\begin{center}I.\\ Den formelle og speculative Logik.\end{center}
\begin{center}II.\\ Den propædeutiske Logiks Formaal.\end{center}

\begin{center}{\bfseries Første Deel.}\\ {\bfseries\emph{Læren om det subjective Begreb.}}\end{center}

\begin{center}{\bfseries\emph{Første Capitel.}}\\ Det subjective Begrebs Dannelse.\end{center}
\S\,1.\quad Almeenforestillingen.\par
\S\,2.\quad Den klare og tydelige Anskuelse.\par
\S\,3.\quad Indsigten eller det subjective Begreb.\par
\S\,4.\quad Begrebets Aprioritet.\par
\S\,5.\quad Begrebets Totalitet.\par

\begin{center}{\bfseries\emph{Andet Capitel.}}\\ Det subjective Begrebs Fremstilling: den logiske Dom.\end{center}
\S\,6.\quad Dommens logiske Form.\par
\S\,7.\quad Dommens logiske Indhold.\par
\S\,8.\quad Dommens logiske Totalitet.\par

\begin{center}{\bfseries\emph{Tredie Capitel.}}\\ Det subjective Begrebs Gyldighed: Slutning og Beviis.\end{center}
\S\,9.\quad Slutningens logiske Reqvisiter.\par
\S\,10.\quad Slutningens formelle Gyldighed.\par
\S\,11.\quad Slutningens reelle Gyldighed: Beviis og Methode.\par

% ---- PDF 9 ----
\begin{center}{\bfseries Anden Deel.}\\ {\bfseries\emph{Læren om det objective Begreb.}}\end{center}
\begin{center}\emph{Indledning.}\end{center}
\S\,12.\quad Den dialektisk-speculative Methode.\par

\begin{center}{\bfseries\emph{Første Capitel.}}\\ Det objective Begrebs Dannelse.\end{center}
\S\,13.\quad Væren og Væsen.\par
\S\,14.\quad Væsen og Virkelighed.\par
\S\,15.\quad Virkelighed og Begreb.\par

\begin{center}{\bfseries\emph{Andet Capitel.}}\\ Det objective Begrebs Frihed.\end{center}
\S\,16.\quad Idealitet og Realitet.\par
\S\,17.\quad Subjectivitet og Objectivitet.\par
\S\,18.\quad Teleologie.\par

\begin{center}{\bfseries\emph{Tredie Capitel.}}\\ Det objective Begrebs Fylde: Ideens Logik.\end{center}
\S\,19.\quad Ideens Oprindelighed.\par
\S\,20.\quad Ideernes Vexelvirkning.\par
\S\,21.\quad Idee og Gud.\par

}
\begin{center}\rule{0.35\linewidth}{0.4pt}\end{center}

\cleardoublepage
\tableofcontents

\mainmatter    % sets arabic numbering itself — do not add \pagenumbering here

% ============================================================
%  BODY — printed pp. 1--283 (PDF 10--292)
% ============================================================

% ---- Indledning, printed pp. 1--5 (PDF 10--14) ----
\division{Indledning.}

% --- p. 1 ---
% The "Indledning." head and its short rule are already in transcription.tex
% above this point. What follows is the head of section I as the page sets it.
%
% IMAGE EVIDENCE for the I./II. heads (printed pp. 1 and 4): a centred bold
% roman numeral ("I." / "II.") in ANTIQUA, then, one line below, the argument in
% BOLD Fraktur, a size up from the body and NOT letterspaced. So these two heads
% sit between \capitel (whose label IS letterspaced and whose argument is NOT
% bold) and \parag. They are numbered I./II., not as §§ -- that asymmetry is in
% the book -- so they are hand-rolled here rather than pushed through \capitel.
\romsec{I.}{Den formelle og speculative Logik.}

% The paragraph opens flush left with a raised initial "V" about 1.5x the body
% size (the same treatment as the "L" of "Lighed" at § 1, printed p. 6). Set as
% ordinary text here; only the absent indent is reproduced.
\noindent Ved at betænke sin egen Tænkning føres Subjectet til at
tydeliggjøre sig de forskjellige Tankeoperationer og disses indbyrdes
Sammenhæng. Den første Fordring er, at Tænkningen skal være conseqvent og
følge sine egne Love uden alt Hensyn enten til det tænkende Subjects
psychiske Eiendommelighed eller til Beskaffenheden af de Gjenstande, paa
hvilke den ellers maatte kunne anvendes. Den formelle Logik behandler derfor
Tænkningen som et givet Aandsphænomen uden at spørge om dens Oprindelse eller
objective Gyldighed. Dette Spørgsmaal er ikke desto mindre saa
betydningsfuldt, at Philosophiens hele Retning og Skjæbne er afhængig af den
% PRINTER'S DEFECT p.1: the final "e" of "Skjæbne" is broken/under-inked and
% prints as two fragments resembling a colon. The word is Skjæbne; transcribed
% so, since this is an inking failure and not a compositor's reading.
Maade, hvorpaa det besvares. Sensualisterne lære, at Sjælen (\textit{tabula
rasa}) faaer sit oprindelige Indhold fra Omverdenens sandselige Indtryk
(Impressioner), og ende med at gjøre Tænkningen til en Afskygning af
Erfaringen. Den subjective Idealisme vindicerer derimod Sjælens uendelige
Selvvirksomhed, og tænker sig Erfaringsverdenen som det forestillende Jegs
eget ubevidste Product. Enhver af disse Theorier miskjen%
% --- p. 2 ---
der den menneskelige Tanke. Sensualismen undervurderer Tænkningen, betvivler
Tankelovenes almene Gyldighed, og kan desaarsag ingen Sandhed finde i den
virkelige Verden; Idealismen overvurderer den subjective Tænkning, benegter
Verdens sande Virkelighed, og ender i formel Tomhed. Begge Eensidigheder
overvindes i den speculative Logik, der viser, at vor fornuftige Tænkning
afspeiler den i Universet aabenbarede Fornuft.

\anm{1} Ved at gjøre Sandsningen til Erkjendelsens Kilde maatte Sensualismen
føre til Skepticisme; thi Sandsningen refererer sig blot til det Particulære
og Tilfældige, ikke til det Almindelige og Nødvendige; Erfaringen lærer blot,
\emph{at} Noget skeer, men ikke, \emph{hvorfor} det skeer; den viser, at det
Ene følger paa det Andet, men beviser derimod ikke, at det Ene er en Følge af
det Andet; under slige Forudsætninger var det altsaa conseqvent, at Hume
bestred Causalitetsbegrebet.

\anm{2} Imod Humes Theorie opponerede Kant. Grundtanken i den kantiske
Fornuftcritik er den, at vore rene Sandseanskuelser (Tid og Rum),
Forstandskategorier (Eenhed og Mangfoldighed, Aarsag og Virkning etc.) og
Fornuftideer (Sjælen, Verden, Gud) ikke kunne være abstraherede af en udenfor
Bevidstheden liggende Erfaringsverden, eftersom de selv udgjøre de nødvendige
Forudsætninger, uden hvilke ingen Erfaring er mulig. Men idet Kant forsvarede
de aprioriske Formers subjective Almeengyldighed, frakjendte han dem paa den
anden Side al objectiv Betydning. Den hele Erfaringsverden bliver følgelig at
betragte som et Product af tvende Factorer, nemlig af de aprioriske Former og
et ubekjendt Object („das Ding an sich“), hvis rene Væsen Ingen kan erkjende,
fordi Enhver, der prøver paa at erkjende det, derved strax optager det i sin
egen Bevidstheds subjective Medium.

% --- p. 3 ---
\anm{3} Det kantiske „Ding an sich“ blev bestridt af den ældre Fichte. Det
er, ifølge Fichtes Anskuelse, en Urimelighed, at antage en „Ding an sich,“ da
en saadan hverken kan sandses eller tænkes. Et Object, der skulde ligge
aldeles udenfor Bevidstheden, kan Subjectet ikke tillægge Realitet uden at
komme i Modsigelse med sig selv; thi hvis Nogen siger, at han veed om et
saadant Objects Tilværelse, da har han allerede draget denne Tilværelse ind i
Bevidsthedskredsen; han siger da: Jeg veed Noget om en Tilværelse, om hvilken
jeg dog aldeles Intet veed. „Das Ding an sich“ var den sidste Rest af
Objectivitet, som den kritiske Idealisme lod blive tilbage. Ved at bortskaffe
denne ubekjendte Størrelse blev Idealismen aldeles subjectiv, og concentrerede
sig nu i den Sætning, at den hele Verden frembringes af Jeget. For at forstaae
denne dristige Paastand, maae vi nærmere bestemme Jegets Væsen, dets
Productivitet og dets Frihed.

% a)--c) are set with a hanging indent: the label stands to the left and the
% continuation lines are indented to the label's text. Not a display head, so
% a plain itemize (no enumitem in the preamble) rather than a \begin{center}.
\begin{itemize}
\item[a)] Det Jeg, hvorom Idealismen taler, er ikke det empiriske Jeg, der
  finder sig umiddelbart bestemt ved Bevidsthedslivets forskjellige Yttringer
  og Tilstande (Jeg taler, sørger, glædes etc.); men det rene Jeg = Jeg,
  Jeget, der i Alt tænker sig selv uden Hensyn til particulære
  Indskrænkninger og Forandringer.
\item[b)] Den Productivitet, hvorved Jeget frembringer den hele virkelige
  Omverden, er intet løst Phantasiespil, eller vilkaarlig Imagination, som
  Jeget efter Behag skulde kunne forandre eller standse, men en
  Nødvendighedsact, hvorved det forudsætter de Skranker, hvis uendelige
  Ophævelse betinger dets Frihed.
\item[c)] Jegets Verdensfrembringelse foregaaer ubevidst i dets egen dunkle
  Livsgrund, og hører, saa at sige, til dets natlige Sphære; derfor træder
  Universet op med en saa
% --- p. 4 ---
  imponerende Magt, at Subjectet tager sit eget Verdenssyn for en udvortes
  gjenstændig Virkelighed. Denne populære Anskuelse er imidlertid kun en
  Drøm, der forsvinder, naar den philosophiske Tænkning vaagner. Overgangen
  fra Forestillingens til Philosophiens Standpunkt er altsaa en Overgang fra
  Illusion til Sandhed, fra Ufrihed til Frihed.
\end{itemize}

\anm{4} Den speculative Logik er i den nyere Tid bleven grundlagt af Hegel.

\romsec{II.}{Den propædeutiske Logiks Formaal.}

% Unlike § I, this section opens with an ordinary indented paragraph: no raised
% initial. Verified on the image.
Det er altsaa en Kjendsgjerning, at der foruden den saakaldte forstandige
Tænkning, der finder sin umiddelbare Anvendelse i det praktiske Liv og sin
høiere Uddannelse i Mathematiken og de empiriske Videnskaber, endnu gives en
anden Art af Tænkning, den dialektiske nemlig, der betinger al høiere Viden,
og derfor nærmest hører hjemme i den speculative Philosophie. Uagtet begge
disse Tænkningens Hovedretninger have deres Udspring fra een og samme Kilde,
den evige Fornuft, danne de dog for den overfladiske Betragtning en saa
paafaldende Modsætning, at det ofte synes, som om Sandheden var i Strid med
sig selv. Denne Modsætning lader sig tilsyne ikke blot i en afvigende
Terminologie, men endog i den Maade, paa hvilken Tankelovene selv bestemmes og
fremstilles, saa at En kan være meget vel øvet i forstandig Tænkning, uden dog
at see sig istand til enten at følge Philosophiens Gang eller tilegne sig dens
Resultater. Skal altsaa den propædeutiske Logik
% --- p. 5 ---
svare til sin sande Bestemmelse, skal den orientere os i Tænkningens hele
Rige, da maa den nødvendigviis indlade sig med begge de modsatte
Tankeretninger. Ved at holde sig til den formelle Tænkning alene vilde den
aabenbart give for Lidet, ved at tabe sig i speculative Problemers Løsning
vilde den overskride sin propædeutiske Grændse: den indtager derfor en
mæglende Plads imellem den formelle og speculative Logik.

\anm{1} For at løse sin Opgave maa den propædeutiske Logik give en comparativ
Fremstilling af de formelle og speculative Tankeoperationer. Skal denne
Fremstilling lykkes, bør det Skridt for Skridt eftervises, hvorledes et
bestemt Tankeindhold, efter først at være betragtet med et forstandigt Øie,
forandrer sit Udseende ved at besees igjennem Dialektikens Glar. \emph{Læren
om det subjective Begreb}, der udgjør Logikens første Deel, lader derfor
Tænkningen trinviis søge sine Udgangspunkter i Sandseverdenens og
Forestillingens fastere Jordbund, for derfra at hæve sig til den speculative
Erkjendelses friere Luftstrøm. Ved denne Bevægelse nedenfra opad føres vi
overalt fra Empirien ind i Philosophien, fra den Tankekreds, vi kalde vor
egen, ind i den Fornuftverden, der aabenbarer den guddommelige Tænker. Den
Tankeøvelse og Indsigt i Begrebets Væsen, der herved vindes, introducerer
dernæst en mere methodisk Udvikling af Tænkningens og Tilværelsens
Universalformer og Principer, der fremstilles i Logikens anden Deel, \dsi{}
\emph{i Læren om det objective Begreb}. Men ligesom den første Hovedafdeling
stedse peger hen til høiere Begrebsudvikling, saaledes maa denne igien overalt
vise tilbage til Begrebernes empiriske Existens, og kaste Lys over den
virkelige Verden.

% Printed p. 5 closes the Indledning with a short centred rule, heavier than
% the rule under the Indledning head. Reproduced.
\begin{center}\rule{0.30\linewidth}{1.0pt}\end{center}

% ---- Første Deel, printed pp. 6--96 (PDF 15--105) ----
% p.6 opens with the double rule + Deel head + Første Capitel + § 1.
% Batch 1 places these; the \deel and \capitel macros are defined in the preamble.

% --- p. 6 ---
% IMAGE EVIDENCE for the p. 6 head, top to bottom: a full-measure DOUBLE rule
% (a thick rule with a thin one below it); "Den propædeutiske Logiks" (large
% bold Fraktur, not letterspaced); "Første Deel:" (bold Fraktur, WITH A COLON,
% NOT letterspaced); "Læren om det subjective Begreb." (largest, bold Fraktur,
% clearly LETTERSPACED); a short centred rule; "Første Capitel." (bold Fraktur,
% letterspaced); "Det subjective Begrebs Dannelse." (one size down, NOT bold,
% not letterspaced); "§ 1." (bold, antiqua figure); "Almeenforestillingen."
% (bold Fraktur, NOT letterspaced).
% The Indhold prints "Første Deel." with a period; the body head prints a colon.
% Both readings are on record in the file header; the colon is what p. 6 sets.
\deel{Første Deel:}{Læren om det subjective Begreb.}
\capitel{Første Capitel.}{Det subjective Begrebs Dannelse.}
\parag{1}{Almeenforestillingen.}

% Flush left, with a raised initial "L" about 1.5x the body size — the same
% treatment as the "V" of "Ved" at the opening of the Indledning, printed p. 1.
% Set as ordinary text; only the absent indent is reproduced.
\noindent Lighed og Ulighed ere de simpleste Former, i hvilke Omverdenens
Gjenstande lade sig fremstille for den tænkende Betragtning. Tankeoperationen,
der bestaaer i at fremhæve og fastholde det Lige i de ulige Gjenstande, er en
\emph{Abstraction}, og det Resultat, til hvilket den fører, en
\emph{Almeenforestilling}, der kan characteriseres som et Sandsebillede, et
Tankebillede og Begrebets Forbillede.

Uagtet tvende tilværende Ting ikke i een og samme Henseende kunne indbyrdes
være baade lige og ulige, saa er det dog umuligt at tænke Ligheden uden ved
Hjælp af Uligheden; jo skarpere disse Kategorier adskilles, desto klarere
viser det sig, at de gaae over i hinanden: deres Grændse er følgelig
dialektisk.

\anm{1} Flere Gjenstande, der have en væsentlig Lighed med hverandre, reducere
vi i Sproget til een Fælles%
% --- p. 7 ---
benævnelse, og vise derved, at vi oversee deres tilfældige Ulighed. Med
Sproget begynder strax Abstractionen; idet Barnet lærer at tale, lærer det
tillige at abstrahere. Gjenstandenes Lighed observeres ved en Sammenligning,
der fra først af er ligesaa umiddelbar som Sandsningen selv, og Abstractionen
foregaaer forsaavidt ganske ubevidst; men da Ligheden selv ikke er nogen
sandselig Gjenstand, der ved udvortes Adskillelse kan fjernes fra de
Uligheder, i hvilke den finder sit Udtryk, saa viser Abstractionen sig som en
ideel Act, hvorved Subjectet selv tager det Almene for det Væsentlige, og
lader de tilfældige Modificationer fare; altsaa som en begyndende
Tænkningsact, der blot behøver at bringes til Bevidsthed, for at den egentlige
Tænkning kan begynde.

\anm{2} At see en Gjenstand er med andre Ord at danne sig en til bestemte
sandselige Indtryk svarende Forestilling. At Subjectet selv danner sig
Forestillingen, indlyser deraf, at Gjenstandens Billede kan reproduceres,
efterat Gjenstanden selv er forsvunden. Ved at sandse en Række af forskjellige
Gjenstande uddanner sig i Bevidstheden en Kreds af forskjellige Billeder og
Forestillinger; ved Anvendelse af Sammenligning og Abstraction erholdes en
Almeenforestilling. Almeenforestillingen er i Begyndelsen ganske umiddelbar.
Man forestiller sig det Almindelige ved at sætte en enkelt Gjenstand som
Exemplar. Almindeligheden falder altsaa sammen med det Billede, i hvilket den
anskues, og er følgelig at betragte som et \emph{Sandsebillede}. Men en
nærmere Undersøgelse viser, at ethvert Exemplar, der skal repræsentere den
rene Almindelighed, baade giver os for Meget og for Lidet. Exemplaret
udtrykker foruden det Almene tillige en bestemt Gjenstands Eiendommelighed, og
indeholder altsaa for Meget; paa den anden Side viser det kun Almindeligheden
i en
% --- p. 8 ---
% PRINTER'S DEFECT p.8 l.1: no full stop after "for Lidet". The sentence space
% before "Subjectet" is a wide one, but the period simply did not print
% (verified at 300 dpi: clean paper, no broken sort). Transcribed as printed.
particulær Skikkelse, og giver os desaarsag for Lidet Subjectet, der nu
forsøger paa at gjøre sig en bestemt Forestilling om det Almindelige, faaer da
kun en uoverskuelig Række af Exemplarer, der, hvert for sig, ere lige adæqvate
og lige uadæqvate. Følgen heraf er, at den med Sandseligheden behæftede
Almeenforestilling forsvinder som et Gjøglebillede, og kun Tanken om det
Almindelige, der selv er en Tanke, bliver tilbage. Ved at gaae over i
Tænkningen bliver Almeenbilledet egentlig ubilledligt. Tænkningen er
forsaavidt negativ, som det Almene tillige er det Usynlige, der ikke lader sig
forestille eller fastholde i noget Billede. Den abstraherende Tænkning frigjør
det Almindelige fra alle sandselige Tilsætninger; men naar det nu endvidere
betænkes, at det Almene aldrig existerer for sig selv, men abstraheres af de
existerende Gjenstande, saa maa Tænkningen atter ophæve Abstractionen, og
Forstandens Øie skue det Almindelige i dets enkelte Exemplarer og Exempler.
Den frie Tanke svæver uafbrudt imellem det Ubilledlige og det Billedlige,
hvorved Almeenbilledet forvandles til et \emph{Tankebillede}. Ifølge heraf
% "omslutter" is NOT letterspaced on the page (zoom-verified) although the
% sense makes it part of the emphasised phrase; the letterspacing begins with
% "og" on the following line. Transcribed as printed.
omslutter \emph{og indbegriber} det alle de enkelte Gjenstande, og er i denne
Betydning at sætte som \emph{Begrebets Forbillede}.

\anm{3} Naar det Almindelige fastholdes ved Abstractionen, kaldes det et
Abstractum; naar det derimod anskues i et bestemt Exemplar, haves det
\textit{in concreto}. Dog bliver Forholdet imellem det Abstracte og Concrete
anderledes bestemt i den formelle end i den speculative Logik.

\anm{4} Hvad Forholdet imellem den rene Lighed og Ulighed angaaer, da er det
for det Første indlysende, at de gjensidig udelukke hinanden, eller at de, med
andre Ord, begrændse hinanden. Men hermed er Sagen dog ikke afgjort; thi
enhver negativ Tankebegrændsning bliver tillige at om%
% --- p. 9 ---
sætte som en positiv Tankebestemmelse. Hvad er det Lige? Sv. Det Lige er det
Ulige uligt, og derved sig selv ligt. Heri ligger da, at Tanken om Lighed ikke
blot falder udenfor, men tillige sammen med Tanken om Ulighed. Den, der direct
tænker paa det Lige, tænker derved indirect paa det Ulige. Uligheden er altsaa
den Kategorie, hvorved Ligheden bestemmes. Grændsen imellem den rene Lighed og
Ulighed kan aldrig staae stille; derfor hedder det i Dialektikens Sprog, at de
gaae over i hinanden. At denne Dialektik ikke er vilkaarlig, beviser
Philosophiens egen Historie. Thi ved at gjøre Grændsen imellem det Lige og
Ulige ubevægelig og udialektisk kunde Eleaterne bevise, at al Vorden var
utænkelig. Skal Vorden tænkes, sagde de, da maa det enten være det Lige, der
kommer af det Ulige, eller det Lige, der fremkommer af det Lige. I sidste
Tilfælde kan ingen Vorden finde Sted, thi da det Lige er = det Lige, saa er
her ingen Overgang; i første Tilfælde maatte Noget blive til af Intet, thi i
det Ulige kan det Lige slet ikke indeholdes. Den eleatiske Abstraction
forraader sin Eensidighed derved, at den ikke vil lade det ene Begreb
integrere det andet.

\parag{2}{Den klare og tydelige Anskuelse.}

Naar Gjenstanden for den tænkende Betragtning viser sig som en \emph{Eenhed} i
en Mangfoldighed, bestaaer den subjective Tankevirksomhed: a) i en
\emph{Analyse} eller ideel Udsondring, hvorved de mangfoldige Kjendemærker og
Bestemmelser, der med Hensyn til Gjenstanden ere at adskille, sættes ud fra
hverandre; b) i en ideel Sammenfatten eller \emph{Synthese}, hvorved den
udsondrede Fleerhed føres tilbage til sin Eenhed. Det ved denne
% --- p. 10 ---
dobbelte Tankeoperation vundne Resultat er en \emph{klar og tydelig
Anskuelse}. Formedelst Anskuelsens Klarhed bliver den gjenstændige Eenhed
udsondret af Tilværelsens uendelige Mangfoldighed og sat som bestemt
Enkelthed; ved Anskuelsens Tydeliggjørelse bliver den enkelte Gjenstand selv
bestemt ved en Mangfoldighed af Enkeltheder, der indbefattes i dens egen
Eenhed. Uden en, om end ubevidst, Analyse og Synthese er ingen bestemt
Sandsning eller Erfaring mulig; men Tænkningens analytiske og synthetiske
Virksomhed er fra Begyndelsen af saa aldeles fortabt i Sandsningen, at den
klare og tydelige Anskuelse selv bærer Sandselighedens umiddelbare Præg.
Imidlertid er Begrebet dog ogsaa her forbilledlig tilstede; thi Analyse og
Synthese kunne ikke anvendes uden paa Forhold; men Forholdet selv er ingen
sandselig Gjenstand.

Naar de umiddelbare Forholdskategorier: Eenhed og Mangfoldighed, selv gjøres
til Gjenstand for en tænkende Undersøgelse, opstaaer et dialektisk Problem.

\anm{1} Forsaavidt det Almene er den Fællesform, i hvis Lighed de enkelte
Exemplarers indbyrdes Ulighed forsvinder, viser det sig som den Eenhed,
hvorunder Gjenstandene kunne subsumeres; forsaavidt det derimod tillige finder
sit concrete Udtryk i ethvert Exemplar for sig, bliver dets Lighed selv
nærmere bestemt ved Uligheden, holder sig i en Gjentagelse, og gaaer over i
Fleerhed. Derfor have ogsaa de fleste \textit{nomina apellativa} baade en
% "apellativa" with a single p is what the page sets; not corrected.
Singularis og en Pluralis.

\anm{2} Naar det, der opfattes, ikke refererer sig til Synets, men til
Hørelsens Sands, kan der vel ikke være
% --- p. 11 ---
Tale om Anskuelse, men dog altid om en klar og tydelig Opfattelse.

\anm{3} Naar der tales om et Forstandens Øie, eller om et Tankesyn, maae disse
Udtryk ikke blot forstaaes billedligt; thi Overgangen fra Sandsning til
Tænkning kan ikke skee ved et Spring. Som der ligger en Tænkning til Grund for
al Sandsning, saa ligger der ogsaa en ideel Sandsning til Grund for den frie
Tænkning; som der gives en sandselig Anskuelse, saa gives der ogsaa en
intellectuel Anskuelse, og ligesaavel som der kan være Tale om en physisk
Blindhed, kan der ogsaa være Tale om en mental Blindhed. Men naar det nu med
Bestemthed skal udsiges, hvad det Ideelle er, mod hvilket Tankens Blik vender
sig, saa møder strax den Vanskelighed, at Tanken aldrig dvæler ved Noget, der
kan henstilles og paapeges som et bestemt Exempel eller en hvilende Skikkelse.
Naar Sandsen slaaer sig til Ro ved en begrændset Gjenstand, saa er Tanken
derimod i evig Uro, og overskrider enhver Grændse. Idet det sandsende Øie
dvæler ved den ene Ting, har det med det Samme vendt sig bort fra den anden;
naar Tanken derimod skal fastholde en Bestemmelse, tager den altid en anden
til Hjælp. Derfor er det klare Tankesyn egentlig et Dobbeltsyn: et
\emph{Hensyn}. Hvad \emph{Hensynet} er fra det subjective Synspunkt, er
\emph{Forholdet} fra det objective; derfor gaaer Tænkningen stedse ud paa at
finde de bestemte Forhold. Som Følge heraf faae Analyse og Synthese en ganske
anden Betydning for den intellectuelle end for den sandselige Anskuelse. For
Sandsen er Anskuelsen klar, naar Gjenstanden i dens rene Enkelthed træder frem
med bestemte Omrids; men Tænkningen glemmer derimod ikke, at den udsondrede
Gjenstand selv er een iblandt flere, at dens Enkelthed altsaa kun kan
bestemmes med Hensyn til andre
% --- p. 12 ---
Enkeltheder, og at den derfor maa sættes som en \emph{særskilt} Gjenstand.
Fra den sandselige Side er Anskuelsen tydeliggjort, naar den i Gjenstandens
Eenhed indeholdte Fleerhed træder frem som en Sum af Enkeltheder; for Tanken
derimod kan Anskuelsen først tydeliggjøres derved, at de givne Enkeltheder
bestemmes med Hensyn til hverandre, og sættes som Led i særskilte Forhold.

\anm{4} Eenhedens Forhold til Fleerheden bestemmes ved den blot forestillende
Tænkning saaledes, at Noget, der i een Henseende kan betragtes som en
Fleerhed, kan i en anden Henseende sættes som en Eenhed; men den dialektiske
Tænkning nøies ikke med en saa udvortes Sammenføining af modsatte Hensyn; den
spørger, om ikke Eenheden, naar den klart og tydeligt gjennemtænkes, selv
bestemmer sig som en Fleerhed. I Platos Dialog, Parmenides, har Socrates
truffet Problemet, hvilket han fremsætter omtrent saaledes: At jeg, Socrates,
i een Henseende er en Eenhed, i en anden Henseende derimod en Mangfoldighed,
er ikke saa forunderligt; men kunde Nogen derimod bevise mig, at ἕν var πολλὰ
og πολλὰ ἕν, da vilde jeg høiligen forundre mig (ἀγαίμην ἂν θαυμαστως, ὦ
% GREEK, zoom-verified at 300 dpi: ἕν carries a rough breathing and an ACUTE
% (the stroke leans the opposite way from the grave of πολλὰ, and the same way
% as the ή of Ζήνων); ἂν carries a breathing and a GRAVE. "θαυμαστως" is
% printed with a plain ω — NO circumflex; transcribed as printed (cf. the
% Rettelser entry for p. 241, θετον -> θεῖον, where the same fount drops an
% accent).
% The fount's script theta ϑ is normalised to θ per the repo's standing
% convention (see texts/nielsen/religionsphilosophie/transcription.tex): it is a
% fount variant, not a distinct letter. This is not optional here — U+03D1 is a
% FATAL error under textalpha/greek-fontenc in text mode.
Ζήνων). Eleaterne (Zeno) efterviste den Modsigelse, der ligger i Fleerhedens
Kategorie. „Hvis det antages, at flere Ting existere, da maa deres Antal baade
være bestemt og ubestemt. Tingene maae tænkes at udgjøre et bestemt Antal, thi
de maae netop være saa mange, som de ere (hverken flere eller færre); men paa
den anden Side maa deres Antal tillige tænkes at være ubestemt; thi imellem de
adskilte Ting gives der endnu andre Ting, og udenfor de bestemte Tings Antal
igjen et Antal af bestemte Ting; følgelig bliver Tingenes Antal uendeligt,
ubestemmeligt og derved i sig selv ubestemt. Da det nu er en Modsigelse, at
tænke Tingenes Antal baade som bestemt
% --- p. 13 ---
og ubestemt, saa antages Fleerheden at være utænkelig“. Sandheden heri er, at
den abstracte Analyse kan gaae i det Uendelige: enhver Fleerhed bestaaer af
Eenheder, enhver af disse Eenheder er igjen at tænke som en Fleerhed med sine
Eenheder o. s. f. Den abstracte Synthese gaaer ligeledes i det Uendelige; thi
den Eenhed, i hvilken en Fleerhed samles, er igjen et Led i en Række af flere
lignende Eenheder; forbindes nu disse ved en ny Synthese, da haves atter en
Eenhed i en ny Række o. s. v. I denne Henseende maae vi altsaa give Eleaterne
Ret, naar de paastaae, at Tingenes Antal vil blive uendeligt, hvis man
overhovedet antager en Fleerhed. Men Tanken om en saadan Uendelighed er ikke i
nogen uopløselig Modsigelse med den Sætning, at Antallet tillige maa være
bestemt. Thi naar Synthesen overalt forbindes med Analysen, bliver enhver
Eenhed bestemt ved \emph{sin} Fleerhed, og Tingenes Antal som Følge heraf
\emph{uendelig bestemt}.

\anm{5} Betragtes Forholdet imellem Eenhed og Fleerhed reent qvantitativt,
bliver Eenheden sat som Continuitet, Fleerheden derimod som Discretion. Enhver
Linie forestiller en saadan Continuitet, hvis Discretion igjen er at søge i
Liniens utallige Punkter.

\parag{3}{Indsigten eller det subjective Begreb.}

Skal en existerende Gjenstand tillige være Gjenstand for en begribende
Tænkning, maa dens Indhold lade sig opfatte i Form af \emph{Heelhed og Dele}.
Idet Heelheden og dens Dele sees gjensidig at betinge og bestemme hinanden,
fremkaldes derved Spørgsmaal om Forudsætning og Conseqvens, om Grund og Følge,
og
% --- p. 14 ---
Tænkningen, der nu gjør sig selv Regnskab for det hele Indbegreb af bestemte
Forholde, viser sig som en \emph{Reflexion}, hvorved Subjectet kommer til et
\emph{Begreb om Tingen}, eller til en \emph{Indsigt} i Tingenes Væsen og
Beskaffenhed.

Den Modsigelse, der ligger i Totalitetens eget Begreb, blev allerede tidlig
fremhævet af de græske Philosopher, ligesom den ogsaa igjennem hele
Philosophiens Historie spiller en vigtig Rolle, navnlig med Hensyn til
Spørgsmaalet om Materiens uendelige Deelbarhed. Som en Modsigelse imellem det
Enkelte og Sammensatte reduceredes den af Kant til en Antinomie, der igjen
fandt sin dialektiske Opløsning i den hegelske Logik.

\anm{1} Naar Eenheden sættes som Heelhed, er hermed udtrykt, at den omfatter
en Mangfoldighed af Enkeltheder, der udgjøre dens Dele. For den sandselige
Betragtning viser det Hele sig blot som Delenes Aggregat; for den
intellectuelle Opfattelse er det Hele derimod at bestemme som Delenes
\emph{almindelige} Sammenhæng, og Delene som Led i Heelhedens
\emph{særskilte} Sammenhæng. Naar Heelheden er given som Forudsætning, følger
Delenes Beskaffenhed og gjensidige Forhold med nødvendig Conseqvens. Derfor
hedder det i Professor \emph{Sibberns} explicative Logik, at man ved at
erkjende det Heles \emph{Tilværelsesmaade} kommer til at indsee Delenes
\emph{Tilværelsesgrund} (\textit{ratio essendi}). Naar Noget nemlig sættes
% "Grunden" and "Følgen" in the next sentence are NOT letterspaced —
% zoom-verified; spacing.py flags them, but the line is simply loosely justified.
direct saaledes, at et Andet dermed er sat indirect, da kaldes hiint Grunden,
dette Følgen; derfor er Grunden altid Forholdets primitive, og Følgen dets
consecutive Led.

\anm{2} Af det bestemte Forhold imellem Heelhed og Dele indlyser det, at det
egentlige Tankesyn er et Dobbeltsyn. Idet Heelheden sees umiddelbart, komme
Delene middel%
% --- p. 15 ---
bart tilsyne, og omvendt. Delene og det Hele gjenspeile eller reflectere
hinanden; derfor har den frie, selvbevidste Tænkning ogsaa faaet Navn af
\emph{Reflexion}.

\begin{itemize}
\item[a)] Ordet Reflexion er vel nærmest hentet fra Optiken; men det
  naturlige Lys er tillige et Afbillede af Tankens Lys. Der er samme Forhold
  imellem Tænkningens modsatte Led, som imellem „to Speile, af hvilke det ene
  reflecterer det andets Billede, og i denne Reflexion reflecterer sit eget“
  % PRINTER'S DEFECT p.15: "Pg.." — two full stops after "Pg", zoom-verified.
  % Transcribed as printed.
  (cfr. Professor Heiberg: Grundtræk til Philosophiens Philosophie, Pg.. 34).
\item[b)] Reflexionen er tillige Udtrykket for den \emph{overveiende og
  prøvende} Tænkning. Denne Ordets Bemærkelse udledes af den foregaaende. Kun
  derved, at den ene Tanke sees i den andens Lys, indlyser det, hvorvidt de
  kunne bestaae med hinanden eller ikke.
\item[c)] Endelig kan Reflexionen ogsaa betegne den \emph{raisonnerende og
  begrundende} Tænkning; thi naar tvende Tanker stride mod hinanden, da kommer
  det an paa, hvilken af disse Subjectet skal opgive, og hvilken det skal
  fastholde; dette kan alene afgjøres ved et bestemt Hensyn eller ved en
  Forudsætning, der indeholder den bestemmende Grund. Tanken om Menneskets
  Frihed lader sig, til Exempel, ikke forsone med Tanken om en uundgaaelig
  Skjæbne: det tænkende Subject maa nu enten fastholde Friheden eller
  Skjæbnen; tages der Hensyn til Samvittighed og Tilregnelighed, saa maa
  Tanken om en uimodstaaelig Skjæbne opgives; tages der derimod blot Hensyn
  til en mechanisk Naturnødvendighed, maa Tanken om Frihed forsvinde. Hvilket
  Hensyn Subjectet i et givet Tilfælde vil tage, kan igjen beroe paa andre
  Hensyn, og saa fremdeles i det
% --- p. 16 ---
  Uendelige. Da ethvert Hensyn indeholder en partiel Grund, saa fører den i
  udvortes Hensyn fortabte Tænkning til et Raisonnement af Grunde. Forsaavidt
  den virkelige Omverden viser os Tingene i deres mange forskjellige
  Combinationer og Forviklinger, bliver den raisonnerende Tænkning nødvendig i
  alle empiriske Undersøgelser. Men den raisonnerende Tænkning er dog af
  endelig Natur; forsøger den at løsrive sig fra den givne Erfaring, da maa
  % "pro" and "contra" are antiqua; the "og" between them is Fraktur.
  den ende med at anføre Grunde \textit{pro} og \textit{contra}, uden
  nogensinde at naae den høieste og sidste Grund. Først naar alle ydre Hensyn
  forsvinde i eet Grundsyn, antager den reflecterende Tænkning en philosophisk
  Characteer, og finder den uendelige Forudsætning, af hvilken alle endelige
  Forholde lade sig deducere. Fra Endelighedens Synspunkt er Tænkningen
  grundig, naar den opfatter Sagen i dens mange forskjellige, saavel fjernere
  som nærmere, Relationer; fra Uendelighedens Standpunkt er den kun da
  grundig, naar den søger og finder den absolute Grund. Som det er en
  philosophisk Affectation, at ville anvende Dialektikens uendelige Former paa
  Hverdagslivets endelige Indhold, eller forsmaae et Raisonnement af Grunde,
  naar Talen er om en forstandig Kritik eller en Orienteren i praktiske
  Forhold, saa røber det paa den anden Side en videnskabelig Borneerthed, at
  ville hjælpe sig igjennem med en blot raisonnerende Tænkning, naar det
  gjælder en Erkjendelse af de evige Sandheder.
\end{itemize}

\anm{3} Resultatet af den gjennemførte Reflexion er \emph{Begrebet}, som
saadant, eller det egentlige Tankeresultat. Begrebet er altsaa et Tankehele
eller en Grundtanke, der lader sig tilsyne i en Mangfoldighed af særskilte
Tankebestem%
% --- p. 17 ---
melser og disses indbyrdes Sammenhæng. Spørges der nu, hvor Begrebet
existerer, da kan det for det Første siges at have sin Tilværelse i det
tænkende Subjects Bevidsthed. Subjectet seer Gjenstanden i dens forskjellige
Relationer, abstraherer af den de bestemte Forhold, sondrer (Analyse) og
forener (Synthese) de partielle Tanker i deres Grundtanke (Reflexion), og
siges derfor at danne sig et \emph{Begreb om Tingen}. Men forsaavidt det er
den existerende Gjenstand, af hvilken det subjective Begreb abstraheres, kan
der ogsaa tales om Gjenstandens eget Begreb, eller om det i Gjenstanden
udtrykte Tankehele. Forskjellen imellem Gjenstanden og dens Begreb er egentlig
talt en Forskjel imellem dens intellectuelle og dens sandselige Side, der
nærmere er at bestemme som en Modsætning imellem Tingens \emph{Indre} og dens
\emph{Ydre}. Det Ydre er for Sandsen, det Indre for Tanken; derfor siges det
subjective \emph{Begreb om Tingen} ogsaa at være en \emph{Indsigt i Tingen}
(ἐπιστήμη --- νοητόν).

\anm{4} For at danne sig et Begreb om de Phænomener, der stykkeviis komme
tilsyne, uden at Heelheden umiddelbart kan overskues, bliver det ofte
nødvendigt, at imaginere sig en intellectuel Heelhed, eller at opstille en
\emph{Hypothese}. Den Erkjendelse, hvortil Hypothesen fører, er imidlertid kun
endelig og relativ. Tanken svæver imellem Empirie og Reflexion. I empirisk
Henseende er Hypothesen utilstrækkelig; thi Delene (det Ydre) \emph{betinge},
men \emph{begrunde} ikke Heelheden (det Indre), og i intellectuel Henseende
støtter Indsigten sig til et Raisonnement af Grunde, der mangler den
tilstrækkelige Grund. Hypothesen giver derfor kun en Sandsynlighed, der maa
hæves til Vished enten ved en ny Erfaring, der da indleder en ny
Tankeundersøgelse, eller ved en dybere Indsigt i Sammenhængens væsentlige
% NB the sentence, the paragraph and Anm. 4 all run on into printed p. 18.
% When the pp. 18--30 fragment is spliced, the blank line that transcription.tex
% keeps between the two markers must be closed up, or a paragraph break will be
% introduced here that is not in the book.
% --- p. 18 ---
Grund (den copernicanske Hypothese, Newton). Hvad Hypothesen er i Forhold til
Heelheden, er \emph{Conjecturen} i Forhold til et enkelt, i Heelheden
manglende, Led. (Jfr. Sibberns Logik \S\,23.)
% "Forhold til" on this line is NOT letterspaced — the ABBYY layer spaces the
% whole line out, but the line is merely loosely justified; only "Conjecturen"
% is Sperrsatz (zoom-verified at 300 dpi).
% "Jfr." here has a descending Fraktur capital = J; on p.24 the same
% abbreviation is set lowercase "jfr.". Both are as printed.

\anm{5} Den forestillende Tænkning bestemmer Forholdet imellem det Hele og
Delene saaledes, at det Hele sættes = Delenes Sum, men større end enhver af
dets enkelte Dele. Denne Synsmaade har sin fuldkomne Gyldighed, saalænge der
blot handles om endelige Størrelser; men er aldeles utilstrækkelig, saasnart
Undersøgelserne føres over paa Uendelighedens Gebeet. Spiritualisterne
(Berkeley), der negtede Materiens Væren, beraabte sig paa den Modsigelse, der
ligger i Tanken om en uendelig Deelbarhed. Materien kan, ifølge deres
Paastand, slet ikke existere; thi da den tænkes at bestaae af uendelig mange
uendelig smaa Dele, maa enhver Deel blive = det Hele. Sæt t. Ex. \textit{a}
% The inline symbols on this page are antiqua against the Fraktur: the
% placeholders a and A are upright roman (set here as \textit{}, following the
% convention fixed for A=A), while the relation signs are the printer's own
% "<" and the lemniscate for infinity. Those two are set in math mode rather
% than typed as raw Unicode: neither U+221E (INFINITY) nor the less-than sign
% is declared in the preamble, and the raw infinity glyph would be a fatal
% error on the real build. The "=" is the ordinary equals, kept as a literal.
(Delen) $<$ \textit{A} (Heelheden); ved at multiplicere med $\infty$ paa begge
Sider, erholdes $\infty$ \textit{a} = $\infty$ \textit{A}, og ved igjen at
dividere begge Factorer med samme Størrelse, nemlig $\infty$, udkommer:
\textit{a} = \textit{A}. Mathematikernes Indvending, at en saadan Operation er
absurd, kan ikke modbevise Spiritualisternes Theorie, thi denne støtter sig
netop til den indrømmede Absurditet. Ligesaa utilstrækkelig er den almindelige
Bemærkning, at der i det Uendelige selv ikke kan sættes nogen Grændse, at
altsaa Storhed og Lidenhed maae flyde sammen i Uendelighedens Taage; thi
Opgaven er jo netop at lade Tankens Lys adsprede denne Taage, eller, med andre
Ord, at erkjende, hvorledes det Endelige og det Uendelige gjensidig bestemme
hinanden. --- Endnu mere bestemt traadte Problemet frem i den kantiske
Fornuftkritik, hvor det reduceredes til en i Substansens Begreb liggende
Antinomie:
% --- p. 19 ---

\emph{Thesis: Enhver sammensat Substans maa bestaae af usammensatte Dele.}
Hvis ikke, da maatte enhver Deel igjen være sammensat, og Sammensætningens
Opløsning fortsættes saalænge, indtil Substansen tilsidst kom til at bestaae
af Intet.

\emph{Antithesis: Der existerer ingen usammensat Substans.} Thi ingen
Sammensætning er mulig uden i Rummet; gaves der nu en Substans, der bestod af
usammensatte Dele, da maatte Rummet selv bestaae af udeelbare Dele; men Rummet
kan deles i det Uendelige.

Sandheden i den kantiske Antinomie er, som Hegel ogsaa har viist, denne, at
Tanken om et Hele stedse medfører Tanken om Dele. Naar en vis bestemt Deel
adskilles fra sin Heelhed, bliver den selv sat som et Hele. Det Samme, der fra
et Synspunkt maa betragtes som et Hele, maa fra et andet Synspunkt bestemmes
som en Deel; Relationen vender sig uophørlig om, og enhver endelig Grændse
ophæves. Men hvergang et bestemt Hele sættes, træder det i en bestemt
Modsætning til sine Dele, saa at Grændsens endelige Side ogsaa kommer tilsyne.
Alt beroer da kun paa, at Forholdet selv nærmere bestemmes. Forholdet imellem
Heelhed og Dele har en vis Lighed med Forholdet imellem Eenhed og Fleerhed,
Continuitet og Discretion; ikke desto mindre er her dog en væsentlig Forskjel
at observere. Naar nemlig Continuiteten er sat, saa er Discretionen
forsvunden; og naar Discretionen sættes, afbrydes Continuiteten. Continuitet
og Discretion kunne ombyttes med hinanden i det Uendelige, men de reflectere
dog ikke hinanden i egentlig Forstand. Dette er derimod Tilfældet med det Hele
og Delene; naar det Hele er sat, saa ere Delene ogsaa satte, og omvendt.
Derfor kaldes dette Forhold et \emph{Reflexionsforhold} eller et
\emph{væsentligt Forhold}.
% "Derfor kaldes dette Forhold et" is NOT letterspaced, though the ABBYY layer
% spaces it — the last two lines of the page are loosely justified. The "2*" at
% the foot is a signature mark, not a note mark.
% --- p. 20 ---

\anm{6} Heelheden er paa dens laveste Trin et mechanisk Compositum, der viser
os Forholdet fra dets endelige Side. Paa et høiere Trin staaer den organiske
Heelhed; derfor siges der ogsaa at være et uendeligt Forhold imellem Liv og
Organisme. I sin høieste Udvikling bliver Heelheden hævet til en fri
Totalitet, en reflecteret \emph{Enkelthed}, hvor Begrebet aldeles falder
sammen med sin Gjenstand. Saaledes er det hele Jeg tilstede i enhver af dets
selvbevidste Livsytringer, og dog er det en Enkelthed, der ikke kan deles. Som
Sandsningen \emph{begynder} umiddelbart med det Enkelte, saa \emph{ender}
Tænkningen med den reflecterede Enkelthed.
% "hævet til en fri Totalitet" is NOT letterspaced (zoom-verified at 300 dpi
% against the adjacent "Enkelthed", which plainly is), though the ABBYY layer
% spaces it out.

\parag{4}{Begrebets Aprioritet.}
% Printed head agrees with the Indhold: "Begrebets Aprioritet." Bold, not
% letterspaced.

Uagtet det subjective Begreb i sin oprindelige Udvikling er betinget ved
Erfaring, saa reflecterer det dog sine egne væsentlige Bestemmelser med en
saadan Selvstændighed og Conseqvens (Begrebsanalyse), at det ingenlunde kan
være et Product eller tomt Gjenskin af Erfaringen. Ligesaa lidet lader
Erfaringsverdenen sig betragte som et Produkt af det subjective Jeg, uagtet al
Erfaring er betinget ved apriorisk Tænkning. Derfor gjør den formelle Logik
ogsaa en Adskillelse imellem \emph{aprioriske og aposterioriske Begreber}. De
aprioriske Begreber tilhøre ikke blot det tænkende Subject, saa at de skulde
danne en uoverstigelig Kløft imellem den subjective Erkjendelse og den
objective Sandhed, men udtrykke de \emph{almindelige} Grundformer, der
% "Product" and, four lines later, "Produkt" — both spellings as printed.
% Only "almindelige" is letterspaced here: "Grundformer" is set solid, though
% the sense wants the whole phrase (zoom-verified; cf. the same lapse at
% "omslutter", p. 8).
forsaavidt ikke forudsætte speciel Erfaring, som de overalt finde deres
Anvendelse paa de
% --- p. 21 ---
existerende Gjenstande og disses forskjellige Forhold. Paa den anden Side maae
de aposterioriske Begreber, der abstraheres af de i Erfaringen forefundne
Gjenstande, heller ikke tænkes at forholde sig ligegyldige til de subjective
Tankelove, thi derved vilde de ophøre at være sande Begreber; de ere meget
mere at betragte som de \emph{særskilte} Former, i hvilke den aprioriske
Tænkning gjenkjender sit eget eiendommelige Indhold (Normalbegreber).

Den uendelige Vexelvirkning imellem den subjective Tænkning og den objective
Fornuft forudsætter en skabende Frihed, som det absolute Prius, af hvilken al
Tænken og Væren have deres fælles Oprindelse.

% PRINTER'S DEFECT: the page range prints as "Pag. 190  200" — the dash between
% the two numerals never got set. Verified at high magnification and high
% contrast: the gap is clean, there is no under-inked rule. Transcribed as
% printed, with the wide space kept.
\anm{1} Professor Sibbern har i sin explicative Logik (2den Udgave Pag.
190\quad 200) skjelnet imellem den \emph{forfølgende}, den \emph{almindelige}
og den \emph{rene Begrebsanalyse}. Den \emph{forfølgende} Analyse fører de i
Gjenstanden forefundne Begrebsbestemmelser ud over den empiriske Grændse, ved
at tydeliggjøre de i samme liggende Conseqvenser. Den \emph{almindelige}
% PRINTER'S DEFECT: "Bgrebet" for "Begrebet" — a dropped e, zoom-verified.
% Transcribed as printed, not corrected.
Begrebsanalyse finder Bgrebet i det levende Ord, adskiller Ordets oprindelige
Betydning fra den afledede, viser dets Slægtskab med andre Ord (Synonymik),
udfinder derved det søgte Begrebs væsentlige Reqvisiter, og formaaer saaledes
ikke blot at bestemme, men endog at corrigere Sprogbrugen. Den \emph{rene}
Begrebsanalyse tydeliggjør Tænkningens egne Kategorier, og spiller en
Hovedrolle i den objective Logik. I alle disse Arter af Begrebsanalyse viser
det sig, at Tænkningen gaaer ud over det Givne, saa at Sensualisterne have
Uret i at negte dens Selvstændighed.

\anm{2} Den kritiske Idealisme, der vel indrømmer
% --- p. 22 ---
Tænkningen en subjectiv Aprioritet, men dog benegter dens Evne til at erkjende
den objective Sandhed, og til Beviis herfor beraaber sig paa Antinomierne,
finder sin Gjendrivelse i disse Antinomiers dialektiske Opløsning. Derimod kan
den subjective Idealisme, der henfører det hele objective Verdensindhold til
Jegets uendelige Aprioritet, ikke gjendrives ved den Kjendsgjerning, at
Tænkningens Udvikling betinges ved Erfaringen, eftersom det ikke er
Erfaringsverdenens Nødvendighed, men dens Selvstændighed, den benegter.

\anm{3} Det cartesianske Princip: \textit{Cogito, ergo sum}, har fundet sin
conseqvente Udvikling i den fichtiske „Wissenschaftslehre.“ Det absolute
Subject, Jeg, kan tvivle om Alt, kun ikke om dets egen Tænken og Væren. Den
Sætning: Jeg \emph{tænker}, falder umiddelbart sammen med den Sætning: Jeg
\emph{er}. Subjectet kan ikke tænke uden at \emph{være} tænkende, og om det
end betvivler alle enkelte Tankers Gyldighed, kan det dog ikke ved
nogensomhelst Tvivl ophæve Visheden af den almindelige Tænkning, eftersom det
for at betvivle al Tænkning nødvendigviis maa tænke, og følgelig stole paa
% The full stop after "Gyldighed" is badly under-inked — barely a speck at 300
% dpi, invisible to both OCR witnesses — but it is there. Transcribed.
% "Wissenschaftslehre" is a German quotation and so stays in Fraktur (no
% italic), exactly as „Ding an sich“ at p. 3; only the Latin "Cogito, ergo sum"
% is antiqua.
Tænkningen i samme Nu, som det betvivler dens Gyldighed. Hvad der ikke har
denne uimodsigelige Vished, maa falde udenfor den egentlige Videnskab. Da
Videnskaben ikke kan være et Aggregat af løse Forestillinger, da der maa gaae
een Grundsandhed igiennem alle enkelte Sandheder, saa maae ogsaa alle
videnskabelige Sætninger lade sig føre tilbage til een uomstødelig
Grundsætning, og af denne igjen deduceres. Har nu Jeget fundet denne
almindelige Grundsætning i sin egen Tænken og Væren, saa har det med det Samme
fundet Sandhedens Kilde i sig selv, og Philosophien, den rene Videnskab,
maatte da ifølge denne Theorie være istand til at begribe den hele Tilværelse
i Kraft af den øverste Grund%
% --- p. 23 ---
sætning, eller, med andre Ord, construere sin Verden \textit{a priori}.
Sandheden heri er, at den videnskabelige Tænkning ikke kan anerkjende Andet
end det, Tankens Conseqvens medfører, og at den Nødvendighed, hvormed
Sandheden afnøder det tænkende Subject sin Anerkjendelse, er den reneste
Prøvesteen for det Visse og Uvisse. Men den subjective Idealismes Construeren
\textit{a priori} medfører desuagtet en Eensidighed, der paa flere Maader
lægger sig for Dagen (Jacobis Strikkestrømpe). Subjectet, der altid har
Blikket henvendt mod den formelle Jeghed, taber den virkelige Verden og dens
concrete Indhold af Sigte. Enten maa den philosophiske Viden da indskrænkes
til et ringe Antal af abstracte Sandheder, eller ogsaa vil Philosophen, der
indlader sig paa at construere Natur og Historie \textit{a priori}, let kunne
overbevises om at have insinueret empiriske Forudsætninger tvertimod det
opstillede Princip. Mislykkede Forsøg paa at construere \textit{a priori},
hvad der kun lader sig erkjende i Erfaringens aposterioriske Medium, have
fremfor Alt bidraget til at bringe Philosophien i Miscredit.

\anm{4} Adskillelsen imellem aprioriske og aposterioriske Begreber refererer
sig, ifølge den formelle Logik, til deres Kilde eller Oprindelse. Den simpleste
Form, i hvilken en saadan Adskillelse kan fastholdes, er denne, at de
aprioriske Begreber siges at være \emph{forud} for Tingene, paa hvilke de
\emph{anvendes}; de aposterioriske derimod at følge \emph{efter Tingene}, af
hvilke de \emph{abstraheres}, saa at hine have deres Oprindelse fra
Tænkningen selv, disse fra Erfaringen. Men nærmere betragtet kan denne
Modsætning ikke holde sig. Det har allerede viist sig, at Tingens Begreb ikke
er at søge i dens tilfældige Ydre, men i dens Indre \dsi{} i dens intelligible
Væsenhed. At Tænkningen altsaa abstraherer sine empiriske
% Faint pencil underlinings by a former reader run under "forud for" and
% "anvendes" here and under several phrases on pp. 19, 25 and 26. They are not
% printing features.
% --- p. 24 ---
Begreber af Gjenstandene, vil da med andre Ord sige, at den i disses bestemte
Væsen gjenkjender sin egen oprindelige Væsenhed. Erfaringen er altsaa ikke
nogen egentlig Begrebskilde, men kun det Middel, hvorved Subjectets
almindelige Tænkning beriges med et i sig udpræget og concret Tankeindhold.
Begrebet, der med Hensyn til vor Opfattelse kan være aposteriorisk, forsaavidt
en speciel Erfaring er gaaet i Forveien, er dog med Hensyn til Gjenstanden
selv apriorisk, da Tingens Indre og Ydre stedse forholde sig som det Primitive
(\textit{dignitate prius}) til det Consecutive (\textit{dignitate
posterius}). Den aprioriske Vexelvirkning imellem Tingens Begreb og vor
begribende Tænkning kommer tilsyne i Udviklingen af de saakaldte
\emph{Normalbegreber} (jfr. Sibbern \S\,47). Normalbegrebet begynder
\textit{a posteriori}, forsaavidt dets almindelige Grundtræk fra først af
opdages i en empirisk given Gjenstand; men det viser sig tillige fra en
apriorisk Side, forsaavidt det med det Samme indlyser, at den virkelige
Gjenstand ikke svarer til dets Fordringer. Dette Misforhold fører da
Subjectet, der gjennemtænker Begrebets Conseqvenser, til at imaginere eller
construere sig en saa fuldkommen Gjenstand, at det antydede Begreb deri finder
sit adæqvate Udtryk. Den frugtbare Anvendelse, der i Virkeligheden lader sig
gjøre af slige normerende Begreber (t. Ex. i Chrystallographien), beviser
klart, at den grundige Tænkning har fundet Tingenes oprindelige Mening, eller,
hvad der er det samme, at de aprioriske Fordringer, der lede Tænkeren, ogsaa
beherske Virkeligheden. Alle saakaldte aposterioriske Begreber ere, naar de
ikke forvexles med blotte Forestillinger eller sandselige Anskuelser, egentlig
at betragte som specielle Tilsyneladelser af den universelle Fornuft, hvis
Aprioritet er ligesaa objectiv som subjectiv. Med Hensyn til vor Begribnings
Methode
% --- p. 25 ---
kan derimod Adskillelsen imellem aprioriske og aposterioriske Begreber
forsaavidt have sin Betydning, som den minder det tænkende Subject om dets
Endelighed og intellectuelle Begrændsning. Forsaavidt de concrete Begreber
vindes \textit{a posteriori}, maa Erfaringen gaae Haand i Haand med den
begribende Tænkning; forsaavidt de aposterioriske Begreber ere den aprioriske
Tænknings egne Typer, maa Erfaringens sporadiske Indhold lutres i den rene,
universelle Dialektik. Fra den aposterioriske Side er al vor Erkjendelse
stykkeviis, og vor Tænkning discursiv; ifølge Begribningens Aprioritet maa det
Stykkevise igjen ophæves i den videnskabelige Grunderkjendelse. Som der altsaa
maae gives empiriske Videnskaber, der støtte sig til det Aposterioriske, og
følge deres egen selvstændige Methode, saa maa der nødvendigviis ogsaa gives
en Videnskabernes Videnskab, Philosophien nemlig, der, ifølge sin
eiendommelige Methode, lægger an paa at erkjende Sandheden i dens Aprioritet.
Den concrete Gjennemførelse af Philosophien vindes kun igjennem Empirien, og
det Lys, der opklarer Empirien, udstraaler igjen fra Philosophien.

\anm{5} Spørgsmaalet om det subjective Begrebs Forhold til den empiriske
Virkelighed forvandler sig i høieste Instans til et Problem, der gjælder
Forholdet imellem \emph{Tænken og Væren}. Ifølge den subjective Idealisme
% The intervening "og" IS letterspaced here (zoom-verified), unlike the "den"
% and "og" inside the p. 21 run, so this is emphasised as one phrase.
gaaer den reale Væren under i den formelle Tænkning, ifølge den crasse
Empirisme (Sensualismen) gaaer den frie Tænkning til Grunde i den tilfældige
Væren. Da nu enhver af disse Theorier har viist sig i sin Eensidighed; da den
subjective Tænkning ikke kan producere Tilværelsens Fylde, og den objective
Tilværelse kun lader sig tænke, fordi den overalt er underkastet Tankens Love,
men ikke selv kan tænke, efterdi den mangler den subjective Frihed: saa maa
den endelige Modsætning
% --- p. 26 ---
imellem Tænken og Væren føres tilbage til sin oprindelige Grundeenhed. I denne
Urgrund eller oprindelige Grund taber Modsætningen sin Endelighed; thi Tænken
falder paa ethvert Punkt sammen med Væren, og begge ere forsaavidt lige
uendelige. Men paa den anden Side kan denne uendelige Eenhed dog heller ikke
være uden Forskjel. Hvis Tænken og Væren vare i deres dybeste Grund
tautologiske, saa vilde den endelige Adskillelse, der netop skal forklares af
Grunden, være uden al Grund. Resultatet bliver altsaa dette, at den endelige,
\emph{consecutive} Modsætning imellem Tænken og Væren søger sin tilstrækkelige,
\emph{primitive} Grund i deres uendelige, absolute Modsætning \dsi{} i det
Begreb, hvori deres Adskillelse er ligesaa uendelig som deres Forening, saa at
deres ydre Forskjel forvandles til en indre, \emph{levende} Eenhed. Det
absolute Væsen, der formedelst sin frie Tænken producerer den virkelige
Tilværelse, aabenbarer med det Samme sin Tilværelses Realitet i Productet af
sin skaberiske Frihed. I dette Vendepunkt slaaer det \emph{Logiske} over i det
\emph{Theologiske}. Anselms ontologiske Beviis for Guds Tilværelse er en logisk
Fremstilling af det absolute Begrebs uendelige Aprioritet. De Indvendinger,
der til forskjellige Tider (Gaunilo---Kant) ere fremførte mod dette Beviis,
træffe ikke Sagen men Formen. Thi Bevisets adæqvate Form kan egentlig først
gjennemføres i Ontologien eller den objective Logik.

\parag{5}{Begrebets Totalitet.}
% Printed head agrees with the Indhold: "Begrebets Totalitet." Bold, not
% letterspaced. NB § 5 ends the Første Capitel, but the Capitel does NOT close
% within this batch and there is no section-end rule on pp. 18--29: § 5 is
% still running at the foot of p. 29.

Indbegrebet af alle til et Begreb henhørende Bestemmelser udgjør dets indre
Totalitet eller \emph{Indhold}; Antallet af de existerende Gjenstande, i hvilke
Begrebet
% --- p. 27 ---
finder sit Udtryk, bestemmer dets ydre Totalitet eller \emph{Omfang}. Med
Hensyn til Indhold og Omfang har man i den formelle Logik inddeelt Begreberne
i \emph{høiere} (\textit{conceptus superiores}) og \emph{lavere}
(\textit{conceptus inferiores}), saaledes at hine, der have mindre Indhold,
men større Omfang, indbefatte disse under sig. Da denne Inddeling især gjælder
Naturforholdet imellem Slægt (\textit{genus}) og Art (\textit{species}), har
man ogsaa i Almindelighed kaldet de overordnede Begreber
\emph{Slægtsbegreber}, de underordnede \emph{Artsbegreber}.

Det dialektiske Forhold imellem Begrebets indre og ydre Totalitet førte
Middelalderens Scholastikere til en metaphysisk Undersøgelse angaaende
Begrebernes Almindelighed (\textit{universalia}) og Tingenes Enkelthed
(\textit{res, ens, entitas, hæcceitas}). Nominalisterne (\textit{universalia
post rem}) miskjendte Begrebets væsentlige Aprioritet; de eensidige Realister
(\textit{universalia ante rem}) oversaae \textit{principium individuationis}.
Den sande Realisme (\textit{universalia in re}) er tillige den høieste
Idealisme (\textit{universalia in mente divina}). Den guddommelige Tænkning,
der reflecterer det absolute Begreb, er eet med den skabende Almagt, der
frembringer den Virkelighed, i hvilken Begrebets indre Totalitet sondres og
forvandles til en ydre Totalitet af relative Begreber, hvis Existens sættes i
og \emph{med} de existerende Ting.
% Only "med" is letterspaced — the preceding "i og" is set solid
% (zoom-verified), though the sense wants "i og med".

\anm{1} Det virkelige Begreb er et Tankehele, der finder sit Udtryk i en
dertil svarende sandselig Heelhed, det er et Indre, der træder frem i sit
Ydre. Skal nu Forholdet imellem det Indre og Ydre nærmere bestemmes, da træde
disse Hovedsider igjen ud fra hinanden som særskilte Totaliteter. Det Ydre,
der sandses, men ikke tænkes, er ligegyldigt mod
% --- p. 28 ---
% PRINTER'S DEFECT: a stray horizontal bar prints over the "e" of the second
% "det" in the next sentence ("og det sandseligt Reelle") — foul case or a
% piece of loose furniture, not an accent. Transcribed as "det".
det Indre, der tænkes, men ikke sandses; Tingens empiriske eller
haandgribelige Udvortes kan aldrig udledes af dens Begreb: det Intellectuelle
og det sandseligt Reelle ere forsaavidt at betragte som incommensurable
Størrelser. Men paa den anden Side kan Begrebet dog heller ikke træde ud i
Virkeligheden uden at forudsætte et Stof eller Materiale, hvori det realiseres
og kommer tilsyne: ethvert Ydre maa svare til \emph{sit} Indre, og
Gjenstandens empiriske Totalitet tilfredsstille Begrebets Fordringer.
Betragtes Omverdenen med et sandseligt Øie, er den aldeles uden Begreb;
betragtes den derimod med Tankens Øie, da kommer Begrebet tilsyne i de
allerforskjelligste Former. Det uadæqvate Forhold imellem Begreb og
Virkelighed giver sig ikke blot tilkjende derved, at Normalbegreberne hemmes i
deres fulde Fremtræden (t. Ex. Planeternes Perturbationer; den menneskelige
Organismes Mangel paa fuldkommen Sundhed), men viser sig navnlig deri, at
Begrebets Indhold falder, formedelst de endelige Tings ydre Adskillelse, ud
fra sig selv i en saa sporadisk Mangfoldighed af partielle Begreber, at den
subjective Reflexion ikke er istand til at opdage disses concrete Sammenhæng.
Istedetfor at reducere de partielle Begreber til Momenter i det absolute
Begrebs indre Totalitet, lægger den subjective Tænkning an paa at gruppere
efter de forskjellige abstractere eller concretere Grundformer, og derved
tydeliggjøre sig Begrebets ydre Totalitet. \emph{Indsigten} i et Begrebs
væsentlige Bestemmelser (Indhold) lægges til Grund for en \emph{Oversigt} over
de Gjenstandes Antal, der subsumeres derunder (Omfang). Det ligger da i Sagens
Natur, at Indhold og Omfang staae i omvendt Forhold.

\anm{2} Sammenlignes Begreberne indbyrdes med Hensyn til Omfang, da kan det
høiere Begreb ikke siges at
% --- p. 29 ---
indeholde de lavere i sig, ligesaalidt som det udtømmer alle den enkelte Tings
intellectuelle Bestemmelser. De høiere og lavere Begreber adskille sig da
formedelst Gjenstandenes ydre Existens saaledes fra hverandre, at hine som de
abstractere ere ligegyldige mod disse; de forholde sig til hinanden som videre
og snevrere Schemata, saa at den hele Begriben reduceres til en egentlig
Schematiseren.

% PRINTER'S DEFECT: the Sibbern quotation that ends "...har et ganske andet.“"
% has NO opening „ anywhere — verified at high magnification on both the line
% "fører Prof. Sibbern (§ 42), at den intellect..." and the lines before it.
% The closing mark is transcribed as printed, so this fragment carries „1 / “2
% and its quote balance is deliberately -1.
Der gives Tilfælde, i hvilke den subjective Tænkning ikke er istand til at
udlede alle Gjenstandens specielle Begrebsbestemmelser af een
Grundbestemmelse, men maa søge forskjellige Udgangspunkter for Betragtningen,
og lade sig nøie med en Combination af flere sammenstødende Begreber, der hver
for sig tilhøre forskjellige Sphærer. Saaledes anfører Prof. Sibbern
(\S\,42), at den intellectuelle Indtrængen i Chrystallisationsformernes Natur
og Væsen hos et Mineral har sit eget Støttepunkt i den mathematiske
Betragtning, hvorimod dens Undersøgelse i Henseende til Glands, Farve o. s. v.
har et ganske andet.“ Forsaavidt et Grundbegrebs Modificationer hidrøre fra
ydre Sammentræf og empirisk bestemte Forhold, kunne de oversees som
uvæsentlige og ligegyldige; men hvad der saaledes i een Henseende er
uvæsentligt, er fra et andet Synspunkt, under Forudsætning af et andet
Grundbegreb, at betragte som væsentligt, og de forskjellige Hensyn og
Combinationer, den empiriske Tankeforfølgelse kræver, vedkomme ikke Logiken
som saadan.

% Greek zoom-verified at 300 dpi: γένη is epsilon with an acute and no
% breathing; εἴδη has smooth breathing + acute over the iota (ἴ) — i.e. the
% form the Rettelser prescribe for p. 187, already correct here.
Udvikles derimod Begrebernes egne Relationer, da viser det sig, at de høiere
Begreber stedse have deres specifiske Indholdsbestemmelser i de lavere. Dette
bliver allerede udtalt, idet de overordnede Begreber bestemmes som
Slægtsbegreber (γένη), de underordnede som Artsbegreber (εἴδη); thi her
gjælder det, at enhver Slægt maa indeholde sine Arter. Livet, betragtet som
Slægtsbegreb,

% [text to be added: pp. 30--41]

% [text to be added: pp. 42--53]

% [text to be added: pp. 54--65]

% [text to be added: pp. 66--77]

% [text to be added: pp. 78--89]

% [text to be added: pp. 90--96]

% ---- Anden Deel, printed pp. 97--283 (PDF 106--292) ----
% p.97 opens with the double rule + "anden Deel:" (LOWERCASE as printed) +
% Indledning + § 12. Batch boundaries are aligned so that p.97 starts a batch.

% [text to be added: pp. 97--108]

% [text to be added: pp. 109--120]

% [text to be added: pp. 121--132]

% [text to be added: pp. 133--144]

% [text to be added: pp. 145--156]

% [text to be added: pp. 157--168]

% [text to be added: pp. 169--180]

% [text to be added: pp. 181--192]

% [text to be added: pp. 193--204]

% [text to be added: pp. 205--216]

% [text to be added: pp. 217--228]

% [text to be added: pp. 229--240]

% [text to be added: pp. 241--252]

% [text to be added: pp. 253--264]

% [text to be added: pp. 265--276]

% [text to be added: pp. 277--283]

\begin{center}\rule{0.35\linewidth}{0.4pt}\end{center}

% ============================================================
% RETTELSER — PDF 293 (printed [284], unnumbered).
% Reproduced for the record. All nine corrections are APPLIED in the text above,
% each flagged at its site with a % ERRATA comment.
% The p.92 entry misprints "ci" for "ei" in its own right-hand column; that is
% transcribed as printed here.
% ============================================================
\cleardoublepage
\chapter*{Rettelser}
\markboth{Rettelser}{Rettelser}
{\small\setlength{\parindent}{0pt}
Side 80 Linie 3 ovenfra Omsætningen læs: Oversætningen\par
--- \quad 92 \quad --- \quad 13 \quad --- \quad A er ei = A, \quad --- \quad A er ci = --- A\par
--- \quad 93 \quad --- \quad 7 \quad --- \quad Betingelser \quad --- \quad Betingethed\par
--- \quad 112 \quad --- \quad 4 nedenfra \quad Daddel \quad --- \quad Dadel\par
--- \quad 152 \quad --- \quad 8 \quad --- \quad med \quad --- \quad ved\par
--- \quad 154 \quad --- \quad 15 \quad --- \quad Dyr=Rige \quad --- \quad Dyre=Rige\par
--- \quad 187 \quad --- \quad 9 \quad --- \quad εἵδη \quad --- \quad εἴδη\par
--- \quad 206 \quad --- \quad 9 \quad --- \quad en miniatur \quad --- \quad en miniature\par
% The errata page sets both forms with the fount's script theta; normalised to θ
% per the standing convention, as everywhere else in this file.
--- \quad 241 \quad --- \quad 1 \quad --- \quad θετον \quad --- \quad θεῖον\par
}

\end{document}
